\chapter{解决问题}\label{ch:g3:problems}

会做四则运算是一回事；知道故事需要\emph{哪一种}运算，
则是另一回事。本章给出解决应用题的方法——每次都
按同样的四个步骤进行——并教你
如何从表格和象形统计图中读取信息。

\section{四个步骤}

\begin{method}[解决应用题]\label{met:g3:problems:steps}
\begin{enumerate}
  \item \emph{读}题两遍，用自己的话说出
        题目要求什么，已知什么。
  \item \emph{选择}运算：如果不容易判断，
        就快速画一个草图或线段图。
  \item \emph{计算}，可以列竖式，也可以口算，
        并用估算作为一道保险。
  \item \emph{作答}，用完整的句子回答，写上单位——
        并检查答案是否合理（一个小孩不会重
        $300$ kg）。
\end{enumerate}
\end{method}

\begin{example}[哪种运算？]\label{ex:g3:problems:which}
词语提示运算方向，草图则起决定作用：
\begin{itemize}
  \item \emph{合在一起／一共}：加法；
  \item \emph{拿走／多出多少／还缺多少}：
        减法；
  \item \emph{几个几／几倍}：
        乘法；
  \item \emph{平均分／共分几组}：除法。
\end{itemize}
\end{example}

\begin{omfigure}
\begin{tikzpicture}[scale=0.62]
  % addition bar
  \draw[thick, fill=omDef!20] (0,0) rectangle (3.4,0.8);
  \draw[thick, fill=omThm!20] (3.4,0) rectangle (5.8,0.8);
  \node[font=\small] at (1.7,0.4) {$26$};
  \node[font=\small] at (4.6,0.4) {$17$};
  \draw[Stealth-Stealth] (0,-0.35) -- (5.8,-0.35);
  \node[below, font=\small] at (2.9,-0.4) {共？个：$26 + 17$};
  % comparison bar
  \begin{scope}[xshift=9.4cm]
  \draw[thick, fill=omDef!20] (0,0.9) rectangle (5.2,1.7);
  \node[font=\small] at (2.6,1.3) {$52$};
  \draw[thick, fill=omThm!20] (0,0) rectangle (3.5,0.8);
  \node[font=\small] at (1.75,0.4) {$35$};
  \draw[Stealth-Stealth] (3.5,0.4) -- (5.2,0.4);
  \node[below, font=\small] at (2.6,-0.15)
    {多？个：$52 - 35$};
  \end{scope}
\end{tikzpicture}

{\small 线段图能直观地显示运算：并排放在一起的部分，
用加法；两条线段相比较，用减法。}
\end{omfigure}

\begin{example}[两步应用题]\label{ex:g3:problems:twostep}
“玛雅买了 $3$ 包，每包 $6$ 瓶，喝掉了 $2$ 瓶，
还剩多少瓶？”
\begin{enumerate}
  \item 问什么：剩下的瓶数。已知：$3$ 包，每包 $6$ 瓶；
        喝掉 $2$ 瓶。
  \item 画草图：$3$ 组，每组 $6$ 个，再
        拿走 $2$ 个，共两步运算！
  \item 计算：$3 \times 6 = 18$，再算 $18 - 2 = 16$。
  \item 作答：还剩 $16$ 瓶。
\end{enumerate}
算完第一步不要停，要重读题目，
看看故事是否已经讲完。
\end{example}

\begin{example}[没用的数]\label{ex:g3:problems:trap}
“一位面包师今年 $52$ 岁，烤了 $120$ 个面包卷，卖出了 $85$ 个。
还剩多少个面包卷？”年龄只是用来分散你注意力的：
答案只需要 $120 - 85 = 35$ 个面包卷。计算之前，
先整理数据：哪些数是题目真正用到的？
\end{example}

\section{读表格与象形统计图}

\begin{example}[表格]\label{ex:g3:problems:table}
图书馆借出的书：
\begin{center}
\renewcommand{\arraystretch}{1.3}
\begin{tabular}{l|cccc}
 & 漫画 & 小说 & 科普 & 合计 \\
\hline
第 1 周 & $14$ & $9$ & $6$ & $29$ \\
第 2 周 & $11$ & $12$ & $8$ & $31$ \\
\end{tabular}
\end{center}
读表：第 2 周，小说 $12$ 本。由表格计算：两周一共借出
漫画 $14 + 11 = 25$ 本；借书最多的是
第 2 周（$31 > 29$）。
\end{example}

\begin{example}[象形统计图]\label{ex:g3:problems:pictogram}
在象形统计图中，一个图形代表一个固定的数量——先读
图例！

\begin{omfigure}
\begin{tikzpicture}[scale=0.62]
  \node[left, font=\small] at (0,2.4) {星期一};
  \foreach \i in {0,1,2} \fill[omDef] (0.5+\i*0.8,2.4) circle (0.26);
  \node[left, font=\small] at (0,1.2) {星期二};
  \foreach \i in {0,1,2,3,4} \fill[omDef] (0.5+\i*0.8,1.2) circle (0.26);
  \node[left, font=\small] at (0,0) {星期三};
  \foreach \i in {0,1} \fill[omDef] (0.5+\i*0.8,0) circle (0.26);
  \fill[omDef] (0.5+1.6,0) circle (0.26);
  \node[right, font=\small] at (5.4,1.2)
    {图例：一个圆点 $= 4$ 个卖出的蛋糕};
\end{tikzpicture}

{\small 象形图：星期一显示 $3$ 个圆点，所以 $3 \times 4 = 12$
个蛋糕；星期二 $5$ 个圆点，$20$ 个蛋糕；星期三 $3$ 个圆点，$12$
个蛋糕。}
\end{omfigure}
\end{example}

\section{练习}

\begin{exercise}[$\star$]\label{exo:g3:problems:1}
对每个故事，只需\emph{说出}所用的运算（不必计算）：
一群 $48$ 只羊加入另一群 $37$ 只羊；\ $56$ 颗樱桃平均分给
$7$ 个孩子；\ $9$ 盒，每盒 $8$ 支铅笔；\ 阿里有 $71$ 张卡片，
又送出了 $18$ 张。
\end{exercise}

\begin{exercise}[$\star$]\label{exo:g3:problems:2}
用四步法求解：“停车场早上停有 $246$ 辆车，
中午又来 $158$ 辆。现在有多少辆车？”
\end{exercise}

\begin{exercise}[$\star$]\label{exo:g3:problems:3}
求解：“把一根 $90$ m 长的绳子剪成每段 $9$ m 长的
小段，可以剪成多少段？”
\end{exercise}

\begin{exercise}[$\star$]\label{exo:g3:problems:4}
求解：“每张票 $8$ 元。$26$ 张票一共要花多少钱？”
\end{exercise}

\begin{exercise}[$\star$]\label{exo:g3:problems:5}
画线段图，再求解：“艾玛有 $64$ 张贴纸，
雨果比她少 $29$ 张。雨果有多少张？”
\end{exercise}

\begin{exercise}[$\star$]\label{exo:g3:problems:6}
分两步求解：“一箱可装 $45$ 个苹果。杂货商收到 $6$ 箱，
卖出了 $178$ 个苹果。还剩多少个苹果？”
\end{exercise}

\begin{exercise}[$\star$]\label{exo:g3:problems:7}
利用\cref{ex:g3:problems:table}中的表格，两周内一共借出
了多少本科学类图书？第二周比第一周
多借出了多少本小说？
\end{exercise}

\begin{exercise}[$\star$]\label{exo:g3:problems:8}
利用\cref{ex:g3:problems:pictogram}中的象形图，三天
一共卖出了多少个蛋糕？哪一天卖出的蛋糕
最多？
\end{exercise}

\begin{exercise}[$\star$]\label{exo:g3:problems:9}
本题中有一个无关的数据——找出它，再求解：“一本
$210$ 页的书售价 $12$ 元。利奥每天读 $35$ 页。
他需要多少天才能读完这本书？”
\end{exercise}

\begin{exercise}[$\star\star$]\label{exo:g3:problems:10}
“一家人，$2$ 位成人和 $3$ 名儿童，去博物馆参观。成人
票每张 $9$ 元，儿童票每张 $4$ 元。他们用一张 $50$ 元纸币付款，
他们应找回多少钱？”（分三步：成人、儿童，
最后算找零。）
\end{exercise}

\begin{exercise}[$\star\star$]\label{exo:g3:problems:11}
自编一道答案为 $100$ 的两步应用题，用四步法写出
解答过程，并让一位同学来试做。
\end{exercise}
